\subsection{The \bsq{Overview}}\label{sec:OverviewDoc}
The Overview page is not a typical documentation comment; its source is an HTML\index{HTML} document that is referenced through the command line option \mbox{\bsq{\texttt{-overview}}} (see \tqfullvref{sec:Locations} and \autocite{ORACLE_DOC_JAVADOC_MAN:StandardDocletOptions}) when the Javadoc tool\index{javadoc} is called to generate the documentation for your project. It takes the the (fully-qualified) name of that file as the argument; the recommended filename is \bdq{\texttt{overview.html}}.

The file itself should be a valid HTML~5\index{HTML} document; you can omit the \lstinline|<head>| section because the Javadoc\index{javadoc} will copy only a section \lstinline|<main>| into the generated documentation, or, if this does not exist, the whole \lstinline|<body>| section.

Several Javadoc\index{javadoc} tags can be also used in the overview document; for details refer to the section \bdq{Where Tags Can Be Used}in \autocite{ORACLE_DOC_JAVADOC_TAG}.

The content of that overview can be nearly everything, from the project's history through build or installation instructions to the manual on how to use the program or library that is delivered through the project. However, you should not duplicate information that is in the module, package, or class documentation comments. If important, you can set reference links to these.

Since Java~23, the overview document can be also a Markdown\index{Markdown} document (see \autocite{JEP467,ORACLE_DOC_JAVADOC_GUIDE:Markdown}); in that case, the recommended file name would be\linebreak \bdq{\texttt{overview.md}}. 

%==============================================================================
% End of file!!
%==============================================================================

